\subsection{Extending \Frexlib{}.}
\subseclabel{extending frexlib}
The Frex library itself, around \FrexlibLoC{} lines of Idris code,
compiles in around 24 seconds. To evaluate its extensibility, we
assigned an experienced \Frexlib{} developer the task to extend the
library with an involutive monoid frexlet. The development took place
over a period of two weeks, with the code-development phase taking 10
days.

This paragraph is for readers who are interested in the breakdown of
the experiment.  It took the developer around an afternoon to design
the frex data structure and produce a pencil-and-paper proof for
soundness and completeness. However, the developer noticed the
structural simplicity of the frexlet, and conjectured that a more modular
construction might be possible. Within two days, the developer found Jacobs's
axiomatisation of involutive algebras~\cite{jacobs:involution}, and
refactored the pencil-and-paper proof using Jacobs's concepts, though
still specialised to involutive monoids only. Then code development
began, and the developer discovered a new performance bottleneck in
\idris{}, which meant that every new 2-3 lines of \idris{} code took five
minutes or longer to type check. To work around the bottleneck, the
developer proceeded with a three buffer approach, using using separate
buffers for definitions that are completed, currently checked and under construction.
Later, when preparing this manuscript, the developer used
the \ByFrex{} construct to implement the involutive monoid fral from
this frex. This step took an additional half-afternoon, mostly spent on
reducing the construction of the initial involutive
monoid to the initial monoid. The management of notation is cumbersome
and slowed the fral development by perhaps an hour or so.

We view the experiment as noisy due to the effect of the type-checker
bottleneck (now eliminated, as described in the next section), but
draw some tentative conclusions.
%
First, a two-week development time seems acceptable for a shipping
component: the involutive monoid frexlet now forms part of \Frexlib{}.
%
Second, we expect \Frexlib{} will need an overhaul of its
notation-system as it accrues more frexlets. This refactoring might
depend on proposing additional notation-management features to
\idris{} first.
%
Finally, although algebraic generalization delayed the development of
the component, the result of the delay is pleasing: it appears that working
with the generic representations in \Frexlib{} encourages algebraic
thinking that can lead to modular code and even to new theorems.
